\appendix
\section{Proofs and Numerical Validation}
\label{app:proofs}

\subsection{Proof of \cref{prop:expansion}}
Apply Taylor's theorem with Lagrange remainder to
$t\mapsto\loss(f_A(x)+t\,\dAj(x),y)$ on $[0,1]$. There exists
$\xi=\xi(x,y)$ on the segment with
\begin{equation}
 \loss(\fj,y)
 =
 \loss(f_A,y)+\grad\loss(f_A,y)^\top \dAj
 +\tfrac12 \dAj^\top\grad^2\loss(\xi,y)\dAj.
\end{equation}
Rearranging, subtracting $\lambda c_j$, and taking expectations gives
\cref{eq:prop1-exact}. Under the $L$-smoothness assumption,
$|d^\top\grad^2\loss(\xi,y)d|\leq L\|d\|^2$ for every realisation, which
yields \cref{eq:prop1-bound} after taking expectations.

\subsection{Proof of \cref{prop:bregman}}
The Bregman three-point identity states that for any $u,v,w$ in the domain of
$\phi$,
\begin{equation}
 D_\phi(u,v)-D_\phi(u,w)
 =
 \left(\grad\phi(w)-\grad\phi(v)\right)^\top(u-v)-D_\phi(v,w).
\end{equation}
Both sides expand to
$\phi(u)-\phi(v)-\langle\grad\phi(v),u-v\rangle$ minus
$\phi(u)-\phi(w)-\langle\grad\phi(w),u-w\rangle$; all $\phi$ terms cancel.
Set $u=y$, $v=f_A$, $w=\fj=f_A+\dAj$, take expectations, and
subtract $\lambda c_j$ to obtain \cref{eq:prop2}. The final claim follows from
the standard strong-convexity and smoothness bounds on Bregman divergences.

\subsection{Proof of \cref{prop:ce}}
For $\loss(f,y)=-f_y+\log\sum_k e^{f_k}$ we have
$\grad\loss(f,y)=p-e_y$ and
$\grad^2\loss(f,y)=\diag(p)-pp^\top\succeq0$. Since
$(\diag(p)-pp^\top)\mathbf{1}=0$, the quadratic form is invariant to adding a
constant to $v$, so
$v^\top(\diag(p)-pp^\top)v=\Var_p(v)$ for the random variable taking value
$v_k$ with probability $p_k$. Let $M=\max_k v_k$, $m=\min_k v_k$, and
$\mu_p=\sum_k p_k v_k$. The Bhatia--Davis inequality gives
$\Var_p(v)\leq(M-\mu_p)(\mu_p-m)\leq(M-m)^2/4$. Replacing $v$ by
$v-\bar v\mathbf{1}$ (which leaves both the quadratic form and the range bound
valid) makes the vector zero-sum, so $M\geq0\geq m$ and
$\|v\|^2\geq M^2+m^2$. Hence
$(M-m)^2=(M+|m|)^2\leq2(M^2+m^2)\leq2\|v\|^2$ and therefore
$v^\top(\diag(p)-pp^\top)v\leq\tfrac12\|v\|^2$. The bound is attained at
$p=(\tfrac12,\tfrac12,0,\dots)$ and $v=(1,-1,0,\dots)/\sqrt2$, so the constant
$\tfrac12$ is tight. Applying \cref{prop:expansion} with $L=\tfrac12$ and using
$-\grad\loss(f_A,y)=e_y-p_A$ gives \cref{eq:prop3}; the upper bound holds
because the remainder $-\tfrac12 d^\top H d$ is non-positive for $H\succeq0$.

\subsection{Proof of \cref{prop:suff}}
For squared loss the exact identity of \cref{prop:bregman} reads
$\widetilde{\marginal}(j\mid A)
=\loss(f_A,y)-\loss(\fj,y)-\lambda c_j
=2\dAj^\top(y-f_A)-\|\dAj\|^2-\lambda c_j$.
The response $\dAj=\dAj(x)$ is a function of the episode input, so it is
measurable with respect to $(x,A)$ and can be taken outside the conditional
expectation; the tower property gives \cref{eq:prop4} with
$\mu(x)=\E[y\mid x,A]$. Linearity in $\dAj$ and the response-only quadratic
penalty follow immediately. For a general differentiable loss the same
argument leaves $\grad\loss(f_A,y)$ inside the expectation, and this gradient
depends on $y$ beyond the residual, which is why state features are also
required.

\subsection{Numerical validation}
\label{app:validation}
We validate the identities numerically rather than relying on the derivations
alone.

\paragraph{Squared loss (exactness).}
Using the subset-capable nonlinear forecasting expert of
\cref{sec:strong-backbone} on held-out episodes, we compared the exact
marginal $\loss(f_A,y)-\loss(\fj,y)$ with the right-hand side of the
squared-loss identity, and separately verified that the seven response
features of the router reproduce the response statistics. Over 384 held-out
(state, candidate) pairs the maximum absolute deviation is $3.3\times10^{-7}$
in normalised units, i.e.\ floating-point exactness. Unit tests enforce the
same identity and the inertness of unselected candidates.

\paragraph{Cross-entropy (two-sided bound).}
For the frozen in-context classifier of the second task family we collected
92{,}995 (state, candidate) samples together with 2{,}975{,}840 per-query rows
and evaluated the sandwich of \cref{prop:ce}. The bound held in
\emph{every} sample at both aggregation levels (smallest upper slack
$3.0\times10^{-8}$, smallest lower slack $1.0\times10^{-3}$). The first-order
term tracks the true marginal closely: Spearman correlation .971 and
$R^2=.853$ pooled, with per-benchmark Spearman between .970 and .993. The
bound is therefore not merely valid but informative, and the empirical
tightness $\rho=(\marginal-\text{first})/(-\tfrac14\|d\|^2)\in[0,1]$ is
reported in the released artifact.

\paragraph{Response sufficiency.}
For the squared-loss expert we fitted the same ridge head twice on identical
held-out state samples: once with state features only and once with state plus
response features. The response features raise the out-of-sample $R^2$ from
$-0.89$ to $-0.22$ on a held-out episode split (a gain of $0.67$), and the
cross-entropy family shows the same direction with a Spearman correlation of
$.97$ between the first-order term and the realised marginal. Both are the
practical content of \cref{prop:suff}: the response supplies the interaction
term that a state-only score cannot represent.

\begin{table}[t]
 \centering
 \caption{Numerical validation of the response--utility relations. All numbers
 are produced by the released scripts; the cross-entropy check uses the frozen
 in-context classifier of \cref{sec:second-family}.}
 \label{tab:theory-validation}
 \begin{tabular}{lll}
  \toprule
  Check & Object & Result \\
  \midrule
  Squared-loss identity & forecast expert & max.\ deviation $3.3\times10^{-7}$\\
  CE sandwich (\cref{eq:prop3}) & in-context classifier & 0 violations / 92{,}995\\
  First-order vs.\ marginal & in-context classifier & Spearman .971, $R^2$ .853\\
  Response vs.\ state only & forecast expert & $R^2$ gain $+0.67$\\
  Hessian spectral bound & softmax, all $C\le 50$ & maximum $=0.5$\\
  \bottomrule
 \end{tabular}
\end{table}
